% Insert after Section 1; no document environment is included.
% Preamble:
% \usepackage{booktabs,tabularx,tikz,hyperref}
% \usetikzlibrary{arrows.meta}
% Both diagrams are native LaTeX; no external image files are required.
% Checkpoint: main, dbbe536 (30 October 2026).
% Sources: supplied Section 2 draft; current controllers, services, Prisma
% schema and frontend; compose.mvp.yml; frontend proxy configuration.
% Acceptance criteria below describe intended/testable behaviour, not claims
% that every UI scenario has passed. Evidence is distinguished in Sections 3/5.

\section{Updated Requirements and System Design}
\label{sec:design}

The design below reflects implementation checkpoint \texttt{dbbe536}.
Acceptance criteria describe expected behaviour; implementation coverage
and recorded verification are distinguished in Sections~\ref{sec:core}
and~\ref{sec:quality}.

\subsection{Main User Stories and Acceptance Criteria}

\paragraph{Learner.}
As a learner, I want to submit a Writing Task~2 essay and receive feedback.
An eligible learner can save a version-controlled draft and submit
non-empty text within the permitted attempt window and revision quota.
Submission returns HTTP~202 and preserves an immutable revision linked to
the pinned assignment version. Word count is displayed; a 250-word
submission gate is not implemented. Traceability links the Writing editor
to \texttt{Attempt}, \texttt{AttemptDraft} and \texttt{SubmissionRevision}.

\paragraph{Instructor.}
As an assigned instructor, I want to review a submission and confirm an
assessment. The instructor can claim a case, inspect available AI feedback,
enter separate criterion bands and evidence, and confirm a teacher decision
without modifying the AI result. Publication follows the release policy;
stale updates and self-review are rejected. Traceability links the review
page to \texttt{ReviewCase}, \texttt{ReviewDecision},
\texttt{AssessmentResult} and \texttt{AcceptedAssessment}.

\paragraph{Administrator.}
As an administrator, I want to manage account roles and status with
accountability. An authorised change requires a reason, updates the account,
revokes its existing sessions and records an \texttt{AuditEvent}.
Unauthorised changes are rejected and the last active administrator is
protected. Traceability links the administration page to \texttt{User},
\texttt{UserRole} and the audit record.

\subsection{UI/UX Design and Main User Flows}
Role-specific prototypes illustrate governance, learner monitoring and
instructor review; their example totals and scores are illustrative data.
They are design references rather than deployment evidence. The current
learner interface uses top navigation and a collapsible mobile menu;
loading messages, inline alerts and disabled controls communicate state.
Backend checks enforce permissions independently of button visibility.

The learner flow connects catalogue, enrolment, lessons, saved progress,
attempts and feedback. Ordered content display does not imply enforced
sequential completion. The instructor flow connects queue, claim, review,
confirmation and policy-controlled publication. Processing status,
provisional feedback and official results are displayed separately.
Section~\ref{sec:core} identifies the demonstrated flows and remaining gaps.

\subsection{Architecture and Technology Decisions}
The application uses React/TypeScript for role-based interfaces, NestJS for
modular backend services and PostgreSQL/Prisma for relational persistence.
A modular monolith keeps business transactions together; asynchronous
grading can run in a separate worker using the same service code.
PostgreSQL-backed jobs and a transactional outbox provide durable work
tracking without requiring Redis/BullMQ. Provider adapters isolate external
AI services and keep credentials server-side.

\begin{figure}[htbp]
\centering
\begin{tikzpicture}[
  archbox/.style={draw,rounded corners,align=center,text width=2.6cm,
    minimum height=0.85cm,font=\small},
  >={Stealth[length=2mm]}
]
\node[archbox] (client) at (0,1.5) {React client\\API proxy};
\node[archbox] (api) at (3.6,1.5) {NestJS API\\Business modules};
\node[archbox] (db) at (7.2,1.5) {PostgreSQL\\Data, jobs, outbox};
\node[archbox] (worker) at (3.6,-0.2) {Grading worker\\Provider adapter};
\node[archbox] (provider) at (7.2,-0.2) {External AI provider};
\draw[<->] (client) -- node[above,font=\footnotesize] {HTTP} (api);
\draw[<->] (api) -- (db);
\draw[<->] (worker) -- (db);
\draw[<->] (worker) -- node[below,font=\footnotesize] {HTTPS} (provider);
\end{tikzpicture}
\caption{Configured MVP architecture. Public deployment remains unverified.}
\label{fig:system-architecture}
\end{figure}

\subsection{Data Model and Important Invariants}
Figure~\ref{fig:core-data-model} shows the main assessment chain and
enrolment relationships. Additional entities connect review cases and
decisions to teacher results, findings to revision paragraphs, and accepted
skill observations to versioned learner profiles and recommendation runs.

\begin{figure}[htbp]
\centering
\begin{tikzpicture}[
  entity/.style={draw,align=center,text width=2.45cm,minimum height=0.65cm,
    font=\footnotesize},
  >={Stealth[length=1.7mm]}
]
\node[entity] (course) at (0,2.8) {Course};
\node[entity] (assignment) at (3.1,2.8) {Assignment};
\node[entity] (version) at (6.2,2.8) {AssignmentVersion};
\node[entity] (attempt) at (9.3,2.8) {Attempt};
\node[entity] (user) at (0,1.35) {User};
\node[entity] (enrolment) at (3.1,1.35) {Enrollment};
\node[entity] (revision) at (9.3,1.35) {SubmissionRevision};
\node[entity] (evidence) at (3.1,-0.1) {SkillEvidence};
\node[entity] (accepted) at (6.2,-0.1) {AcceptedAssessment};
\node[entity] (result) at (9.3,-0.1) {AssessmentResult};
\draw[->] (course) -- node[above,font=\scriptsize] {1:N} (assignment);
\draw[->] (assignment) -- node[above,font=\scriptsize] {1:N} (version);
\draw[->] (version) -- node[above,font=\scriptsize] {1:N} (attempt);
\draw[->] (attempt) -- node[right,font=\scriptsize] {1:N} (revision);
\draw[->] (revision) -- node[right,font=\scriptsize] {1:N} (result);
\draw[->] (result) -- node[above,font=\scriptsize] {1:0..1} (accepted);
\draw[->] (accepted) -- node[above,font=\scriptsize] {1:N} (evidence);
\draw[->] (user) -- node[above,font=\scriptsize] {1:N} (enrolment);
\draw[->] (course) -- node[above,sloped,font=\scriptsize] {1:N} (enrolment);
\end{tikzpicture}
\caption{Simplified relationship map; arrows run from parent to child.
Other entities and foreign keys are omitted for readability.}
\label{fig:core-data-model}
\end{figure}

Key invariants are immutable submitted revisions, pinned assignment and
grading configurations, and separate AI/teacher assessments. An accepted
assessment is not automatically published. Only current, released evidence
affects the official skill profile, and revisions of one attempt do not
count as independent learning opportunities. Draft/review version checks
and transactional job persistence protect against stale or repeated work.

\subsection{Main APIs and End-to-End Data Flow}
Table~\ref{tab:main-apis} lists direct backend routes. The frontend proxy
adds \path{/api} and removes that prefix before forwarding to NestJS.
IDs are UUIDs; protected requests derive identity from the authenticated
session rather than client-supplied role or learner identifiers.

\begin{table}[htbp]
\centering
\begingroup
\small
\setlength{\tabcolsep}{4pt}
\renewcommand{\arraystretch}{1.08}
\begin{tabularx}{\linewidth}{@{}l p{0.43\linewidth}X@{}}
\toprule
\textbf{Method} & \textbf{Backend path} & \textbf{Purpose and access} \\
\midrule
POST & \path{/auth/login} & Authenticate valid account credentials. \\
GET & \path{/courses} & Public course catalogue. \\
POST & \path{/courses/:id/enroll} & Learner enrolment; course must accept enrolments. \\
PUT & \path{/attempts/:id/draft} & Owner saves draft with \texttt{lockVersion}. \\
POST & \path{/assignments/:id/submissions} & Owner submits using \texttt{attemptId}; returns HTTP~202. \\
GET & \path{/teacher/reviews} & Instructor queue scoped to assigned courses. \\
POST & \path{/reviews/:id/claim}, \path{/reviews/:id/confirm} & Assigned instructor claims/confirms a review case. \\
POST & \path{/assessments/:id/publish} & Instructor publishes an accepted assessment. \\
PATCH & \path{/admin/users/:id} & Administrator updates roles/status with a reason. \\
GET & \path{/me/skills}, \path{/me/recommendations} & Learner profile and eligible activity suggestions. \\
\bottomrule
\end{tabularx}
\endgroup
\caption{Selected API contracts at the interim checkpoint.}
\label{tab:main-apis}
\end{table}

The main data flow is draft/attempt validation, atomic revision and grading
request persistence, worker execution and output validation, teacher
confirmation, publication, then evidence projection and recommendation.
When AI is unavailable or unconfigured, review can proceed manually.
Publication uses \texttt{AcceptedAssessment.id}; confirmation uses
\texttt{ReviewCase.id}. Validation, authentication, permission and
concurrency failures return explicit error responses rather than silently
changing the requested result.
